\section{Results}
\label{sec:results}

\subsection{Float32 baselines}
\label{subsec:baselines}
Table~\ref{tab:baselines} reports the float32 ceilings. Fine-tuning helps the
general-purpose encoders dramatically (RoBERTa $+0.434$, MPNet $+0.426$) but
leaves BGE essentially unchanged ($-0.006$, not significant, paired Wilcoxon
$p=0.49$), because BGE is already pre-trained for retrieval. These scores are
the upper bounds for every compression method.

\begin{table}[t]
  \centering
  \caption{Float32 baselines: pre-trained vs fine-tuned, mean MRR@10 over 13
  NanoBEIR subsets. $\Delta$ is the change from fine-tuning.}
  \label{tab:baselines}
  \begin{tabular}{lccc}
    \toprule
    \textbf{Configuration} & \textbf{BGE} & \textbf{RoBERTa} & \textbf{MPNet} \\
    \midrule
    Pre-trained  & 0.691 & 0.097 & 0.156 \\
    Fine-tuned   & 0.685 & 0.530 & 0.582 \\
    $\Delta$     & $-0.006$ & $+0.434$ & $+0.426$ \\
    \bottomrule
  \end{tabular}
\end{table}

\subsection{Post-hoc compression}
\label{subsec:posthoc}
INT8 at $4\times$ is near-lossless on all backbones (largest drop $-0.015$ on
MPNet; full INT8 and TurboQuant tables in Appendix~\ref{app:posthoc}). At
$32\times$, the four post-hoc methods separate clearly
(Table~\ref{tab:posthoc32}, Figure~\ref{fig:crossmethod}). The ordering
PQ $>$ TurboQuant-1bit $>$ post-hoc binary is stable across backbones, and PQ
with $8$-bit codes is best for BGE and MPNet, recovering all but $0.6$ points of
the BGE float32 ceiling. Post-hoc binary is the weakest method everywhere.

\begin{table}[t]
  \centering
  \caption{Four post-hoc methods at $32\times$, mean MRR@10. Best per backbone
  in bold; ties at displayed precision bolded together.}
  \label{tab:posthoc32}
  \begin{tabular}{lcccc}
    \toprule
    \textbf{Method} & \textbf{Ratio} & \textbf{BGE} & \textbf{RoBERTa} & \textbf{MPNet} \\
    \midrule
    Float32 (fine-tuned)         & $1\times$  & 0.685 & 0.530 & 0.582 \\
    Post-hoc binary              & $32\times$ & 0.646 & 0.501 & 0.534 \\
    TurboQuant, 1-bit            & $32\times$ & 0.666 & \textbf{0.519} & 0.547 \\
    PQ ($M{=}192, n{=}4$)        & $32\times$ & 0.675 & \textbf{0.519} & 0.549 \\
    PQ ($M{=}96, n{=}8$)         & $32\times$ & \textbf{0.679} & 0.514 & \textbf{0.562} \\
    \bottomrule
  \end{tabular}
\end{table}

\begin{figure}[t]
  \centering
  \includegraphics[width=0.8\textwidth]{fig_crossmethod_32x.pdf}
  \caption{Post-hoc methods at $32\times$, mean MRR@10 by backbone. The $y$-axis
  starts at 0.48 to highlight inter-method differences.}
  \label{fig:crossmethod}
\end{figure}

\subsection{Training-time binarization beats post-hoc, and sometimes float32}
\label{subsec:bat}
Both BAT variants beat post-hoc binary on every backbone
(Table~\ref{tab:bat}). The gains are largest on the general-purpose encoders:
BAT-STE improves over post-hoc binary by $6.4$ points on RoBERTa and $4.8$ on
MPNet, versus $1.3$ on BGE. Strikingly, BAT \emph{exceeds the float32 ceiling}
on RoBERTa: BAT-STE reaches 0.565 ($p=0.005$) and BAT-atanh reaches 0.590
($p<10^{-3}$), against a float32 baseline of 0.530. BAT-STE matches float32 on
MPNet and BAT-STE is best on BGE, so the best surrogate is backbone-dependent.

\begin{table}[t]
  \centering
  \caption{Binarization-aware training vs post-hoc binary at $32\times$, mean
  MRR@10. The float32 fine-tune is the ceiling. Best $32\times$ result per
  backbone in bold.}
  \label{tab:bat}
  \begin{tabular}{lcccc}
    \toprule
    \textbf{Method} & \textbf{Ratio} & \textbf{BGE} & \textbf{RoBERTa} & \textbf{MPNet} \\
    \midrule
    Float32 (ceiling)        & $1\times$  & 0.685 & 0.530 & 0.582 \\
    Post-hoc binary          & $32\times$ & 0.646 & 0.501 & 0.534 \\
    BAT-STE                  & $32\times$ & \textbf{0.658} & 0.565 & \textbf{0.582} \\
    BAT-atanh ($\gamma{=}0.1$) & $32\times$ & 0.638 & \textbf{0.590} & 0.580 \\
    \bottomrule
  \end{tabular}
\end{table}

\subsection{Why binarization beats full precision on RoBERTa}
\label{subsec:effdim}
The result is counterintuitive, because binarization removes information. We
analyze it in three steps.

\paragraph{The gain is in ranking, not coverage.}
On RoBERTa, MRR@10 rises from 0.530 to 0.565 (STE) and 0.590 (atanh), both
significant (Wilcoxon $p=0.005$ and $p<0.001$; all tests in
Appendix~\ref{app:sig}). Recall@10 barely moves (0.553 to
0.552 and 0.560) and is not significant ($p=0.57$, $p=0.27$). So the same
relevant documents are retrieved, but the correct answers are ranked higher.

\paragraph{The mechanism is effective dimension.}
Encoders without retrieval pre-training produce anisotropic embeddings that
occupy a narrow cone~\citep{ethayarajh2019contextual}. We measure this with the
effective dimension, or participation ratio~\citep{litwinkumar2017optimal,
gao2017theory},
\begin{equation}
  d_{\mathrm{eff}} = \frac{\left(\sum_i \lambda_i\right)^2}{\sum_i \lambda_i^2},
\end{equation}
where $\lambda_i$ are the 768 eigenvalues of the centered embedding covariance.
$d_{\mathrm{eff}}$ counts how many axes carry variance. Float32 RoBERTa is the
most compressed of the three spaces ($d_{\mathrm{eff}}=85$, or 11.0\% of 768).
Binarization-aware training spreads variance across more axes, raising
$d_{\mathrm{eff}}$ to 25.7\% (STE) and 30.0\% (atanh), and MRR@10 tracks this
increase (Table~\ref{tab:effdim}, Figure~\ref{fig:effdim}). Because ranking
under Hamming and cosine depends on angular separation rather than per-axis
precision, one bit per dimension is enough to realize the benefit.

\begin{table}[t]
  \centering
  \caption{Effective dimension $d_{\mathrm{eff}}$ (participation ratio) with
  retrieval quality, on the pooled NanoBEIR corpus. Binarization-aware training
  raises both $d_{\mathrm{eff}}$ and MRR@10 for RoBERTa.}
  \label{tab:effdim}
  \begin{tabular}{lcccc}
    \toprule
    \textbf{Model} & $d_{\mathrm{eff}}$ & \textbf{\% of 768} & \textbf{MRR@10} & \textbf{Recall@10} \\
    \midrule
    RoBERTa-FT        & 85  & 11.0\% & 0.530 & 0.553 \\
    BGE-FT            & 135 & 17.6\% & 0.685 & 0.656 \\
    MPNet-FT          & 102 & 13.3\% & 0.582 & 0.583 \\
    RoBERTa-BAT-STE   & 197 & 25.7\% & 0.565 & 0.552 \\
    RoBERTa-BAT-atanh & 230 & 30.0\% & \textbf{0.590} & 0.560 \\
    \bottomrule
  \end{tabular}
\end{table}

\begin{figure}[t]
  \centering
  \includegraphics[width=0.78\textwidth]{fig_roberta_effdim.pdf}
  \caption{Cumulative explained variance vs number of principal components. A
  curve that rises quickly indicates a low effective dimension. RoBERTa-FT rises
  fastest; RoBERTa-BAT-atanh rises slowest.}
  \label{fig:effdim}
\end{figure}

\paragraph{It is RoBERTa-specific, and structural.}
BGE already has the highest float32 $d_{\mathrm{eff}}$ (17.6\%), so it has little
anisotropy to correct and binarization only costs it quality. MPNet is in
between and is near-lossless under BAT. RoBERTa has the most to gain because it
starts most compressed, which gives the relative ordering RoBERTa $>$ MPNet $>$
BGE. The gain is not under-training: evaluating intermediate checkpoints every
2{,}000 steps (Figure~\ref{fig:trajectory}), both BAT variants sit above the
float32 fine-tune from the first checkpoint, and the gap does not close with
more training.

\begin{figure}[t]
  \centering
  \includegraphics[width=0.78\textwidth]{fig_training_trajectory.pdf}
  \caption{Mean MRR@10 at intermediate RoBERTa checkpoints. The gap from BAT to
  the float32 fine-tune is present early and does not close with training.}
  \label{fig:trajectory}
\end{figure}

\subsection{Matryoshka and stacked compression}
\label{subsec:stacked}
MRL alone is near-lossless to $d=256$ ($3\times$) on all backbones, then drops
sharply below $d=128$ (Figure~\ref{fig:mrl}; full table in
Appendix~\ref{app:mrl}). Its value is that it composes with
a second compressor. Stacking MRL with a precision-reduction method pushes the
frontier well beyond what either axis reaches alone
(Figure~\ref{fig:stacked}). Three results stand out. First, MRL$+$PQ dominates
the high-compression region: at $\geq 96\times$ it is best on every backbone,
and its margin grows with compression (at $384\times$, by 13.0, 6.2, and 7.5
points on BGE, RoBERTa, and MPNet). Second, MRL$+$BAT-atanh beats MRL$+$BAT-STE
on all backbones and all compression ratios, unlike the standalone case where
atanh wins only on RoBERTa; binarizing with a sign before the MRL loss destroys
more magnitude information than tanh. Third, MRL$+$BAT-atanh crosses the float32
ceiling on RoBERTa at $32\times$ (0.560 vs 0.530), the only backbone for which
any stacked configuration does so.

\begin{figure}[t]
  \centering
  \includegraphics[width=0.78\textwidth]{fig_mrl_truncation.pdf}
  \caption{Mean MRR@10 vs MRL truncation dimension $d$. The point at $d=768$ is
  the float32 baseline.}
  \label{fig:mrl}
\end{figure}

\begin{figure}[t]
  \centering
  \includegraphics[width=\textwidth]{fig_stacked_combinations.pdf}
  \caption{Mean MRR@10 vs compression ratio for the five stacked combinations
  and MRL-only (grey dashed). Hollow markers are pure training-time stacks;
  filled markers are hybrid stacks. The dotted line is the float32 baseline.}
  \label{fig:stacked}
\end{figure}

\subsection{Pareto frontiers and cross-backbone generalization}
\label{subsec:pareto}

\paragraph{The Pareto frontier looks different for each backbone.}
Figure~\ref{fig:pareto} plots every configuration as a per-backbone Pareto
frontier of quality against compression. For BGE, no training-time method reaches
the frontier at all. It runs through TurboQuant at moderate ratios, then PQ and
stacked MRL$+$PQ from $32\times$ out past $1500\times$
(Appendix~\ref{app:pareto}). For RoBERTa and MPNet the picture flips: the best
$32\times$ point is a \emph{training-time} method. BAT-atanh leads on RoBERTa
(0.590, 7.1 points above the best post-hoc option), and BAT-STE leads on MPNet
(0.582, 2.0 points above).

\paragraph{The best method does not transfer across backbones.}
To compare backbones of different quality fairly, we divide each score by its own
backbone's float32 score (Figure~\ref{fig:crossbackbone}). The post-hoc methods
keep a similar fraction of quality on all three backbones, so they generalize.
The training-time methods do not: BAT-STE and BAT-atanh help most on RoBERTa,
less on MPNet, and least on BGE. The practical lesson is simple. A study run only on BGE would pick post-hoc PQ
as the best method at $32\times$. But on general-purpose backbones,
training-time binarization is the better choice at $32\times$. So results from
one backbone do not always transfer to another.

\begin{figure}[t]
  \centering
  \includegraphics[width=\textwidth]{fig_pareto.pdf}
  \caption{Quality vs compression Pareto frontiers for BGE, RoBERTa, and MPNet
  ($x$-axis log scale). Filled markers are Pareto-optimal; hollow markers are
  dominated. Colour encodes method family.}
  \label{fig:pareto}
\end{figure}

\begin{figure}[t]
  \centering
  \includegraphics[width=0.95\textwidth]{fig_cross_backbone.pdf}
  \caption{Quality retention (compressed MRR@10 divided by the same backbone's
  float32 MRR@10) by configuration. The rightmost column shows the
  cross-backbone standard deviation.}
  \label{fig:crossbackbone}
\end{figure}
